\section*{Resumen ejecutivo}\label{sec:resumen}
\addcontentsline{toc}{section}{Resumen ejecutivo}

En la Universidad Técnica Estatal de Quevedo, el proceso de titulación de la carrera de Ingeniería de Software va desde la matrícula en Desarrollo de Titulación I hasta la legalización de los documentos finales. En el camino intervienen el estudiante, la comisión formativa, el director, el docente de Titulación I, el tribunal, la coordinación de carrera y asesoría legal. Buena parte de los trámites se hace todavía de forma manual o por correo institucional, con riesgo de perder información, retrasos en las revisiones y errores en el control de los plazos.

En la asignatura de Aplicaciones Web, tres integrantes del equipo construyeron uteq-export, un prototipo web monolítico con Angular, Spring Boot y PostgreSQL que centraliza ese flujo y que todavía no está en uso en la carrera. Las pruebas con usuarios mostraron sus límites: los estudiantes calificaron su usabilidad con 41,88 en la escala SUS, y los tiempos de respuesta solo se midieron sin usuarios simultáneos.

SIC-UTEQ es la versión distribuida de ese sistema y la usarán los mismos actores. La carga se concentra en las fechas límite, cuando muchos estudiantes suben documentos a la vez. Un monolito bien dimensionado podría atender esa carga en una sola máquina; lo que exige distribuir son tres requisitos: escalar solo la recepción de documentos en los picos, evitar que una falla en la evaluación o en las notificaciones detenga las entregas y desplegar un cambio de reglas sin detener el sistema. Por eso el sistema se separa en cinco microservicios por subdominio, cada uno con su base de datos, detrás de un API Gateway, con REST para lo que espera respuesta y RabbitMQ para los avisos.

Este documento corresponde a la primera entrega. Presenta el contexto del problema, los objetivos, 13 historias de usuario, 9 requisitos no funcionales cuantificados con ISO/IEC 25010, la arquitectura en el modelo C4, el modelo de comunicación, el análisis CAP y PACELC, dos registros de decisiones de arquitectura, los riesgos técnicos y un esqueleto con dos servicios que ya se comunican por REST. En el resto del periodo se construirán el gateway, los demás servicios y los contenedores (E2), las bases de datos por servicio con la replicación que pide el análisis CAP y PACELC (E3), y las pruebas de carga y de fallos que comprueban los requisitos no funcionales antes de la defensa (E4).
